\documentclass[aspectratio=169]{beamer}

\usepackage[utf8]{inputenc}
\usepackage[T1]{fontenc}
\usepackage{amsmath,amssymb,amsfonts,amsthm,bm,mathtools}
\usepackage{physics}
\usepackage{braket}
\usepackage{hyperref}

\usetheme{Madrid}

\title{Gibbs Sampling}
\author{Morten Hjorth-Jensen}
\date{FYS5429/9429 Spring 2026}

\begin{document}

\frame{\titlepage}

\begin{frame}{Goals of These Additional Notes}
These notes develop Gibbs sampling from a mathematical viewpoint, with emphasis on:
\begin{itemize}
    \item conditional distributions,
    \item Markov chains and invariant measures,
    \item detailed balance,
    \item convergence intuition,
    \item and connections to statistical physics and Bayesian inference.
\end{itemize}

\medskip

The aim is to explain not only \emph{how} Gibbs sampling works, but \emph{why} it correctly samples from a target distribution.
\end{frame}

\begin{frame}{The Basic Sampling Problem}
Suppose we want to sample from a probability distribution
\[
\pi(x_1,\dots,x_d)
\]
on a high-dimensional space.

\medskip

Typical situations:
\begin{itemize}
    \item a Boltzmann distribution in statistical mechanics,
    \item a posterior distribution in Bayesian inference,
    \item a latent-variable model in machine learning.
\end{itemize}

In many important cases:
\begin{itemize}
    \item \(\pi\) is known only up to normalization,
    \item direct sampling is difficult,
    \item but conditional distributions are simple.
\end{itemize}

This is exactly the setting where Gibbs sampling becomes useful.
\end{frame}

\begin{frame}{Target Distribution and Normalization}
Often the target distribution has the form
\[
\pi(x)=\frac{1}{Z} f(x),
\]
where
\[
Z=\int f(x)\,dx
\]
or, in the discrete case,
\[
Z=\sum_x f(x).
\]

In statistical physics one often writes
\[
\pi(x)=\frac{1}{Z}e^{-\beta E(x)},
\]
where:
\begin{itemize}
    \item \(E(x)\) is an energy,
    \item \(\beta = 1/(k_B T)\),
    \item \(Z\) is the partition function.
\end{itemize}

The difficulty is that \(Z\) is usually hard to compute explicitly.
\end{frame}

\begin{frame}{Why Conditional Distributions Help}
Even when \(\pi(x_1,\dots,x_d)\) is hard to sample from directly, the conditional distributions
\[
\pi(x_i \mid x_1,\dots,x_{i-1},x_{i+1},\dots,x_d)
\]
may be easy to sample.

\medskip

The core idea of Gibbs sampling is:
\begin{quote}
sample one coordinate at a time from its exact conditional distribution, while keeping the others fixed.
\end{quote}

This creates a Markov chain whose stationary distribution is the desired target \(\pi\).
\end{frame}

\begin{frame}{Conditional Probability}
For two variables \(X,Y\), the conditional distribution is
\[
\pi(x\mid y)=\frac{\pi(x,y)}{\pi_Y(y)},
\qquad
\pi_Y(y)=\int \pi(x,y)\,dx,
\]
provided \(\pi_Y(y)>0\).

In \(d\) dimensions, the conditional distribution for coordinate \(i\) is
\[
\pi(x_i \mid x_{-i})
=
\frac{\pi(x_i,x_{-i})}{\int \pi(x_i',x_{-i})\,dx_i'}.
\]

Here \(x_{-i}\) means all coordinates except \(x_i\).
\end{frame}

\begin{frame}{The Gibbs Sampling Update Rule}
Let the current state be
\[
x^{(t)}=(x_1^{(t)},\dots,x_d^{(t)}).
\]

One full Gibbs sweep updates coordinates sequentially:
\[
x_1^{(t+1)} \sim \pi(x_1 \mid x_2^{(t)},\dots,x_d^{(t)}),
\]
\[
x_2^{(t+1)} \sim \pi(x_2 \mid x_1^{(t+1)},x_3^{(t)},\dots,x_d^{(t)}),
\]
and so on until
\[
x_d^{(t+1)} \sim \pi(x_d \mid x_1^{(t+1)},\dots,x_{d-1}^{(t+1)}).
\]

This defines a Markov chain on the state space.
\end{frame}

\begin{frame}{Two-Variable Example}
Suppose the target is \(\pi(x,y)\).

A Gibbs iteration is:
\[
x^{(t+1)} \sim \pi(x \mid y^{(t)}),
\]
\[
y^{(t+1)} \sim \pi(y \mid x^{(t+1)}).
\]

This is already enough to understand the general structure:
\begin{itemize}
    \item each update is exact conditional sampling,
    \item the resulting process is Markovian,
    \item the full joint distribution is preserved.
\end{itemize}
\end{frame}

\begin{frame}{Transition Kernel for One Coordinate Update}
Fix a coordinate \(i\). A single-coordinate Gibbs update defines a transition kernel
\[
K_i(x,dx')
=
\delta_{x_{-i}}(dx'_{-i})\,\pi(dx'_i\mid x_{-i}).
\]

Interpretation:
\begin{itemize}
    \item all coordinates except \(i\) remain unchanged,
    \item coordinate \(i\) is redrawn from its conditional law.
\end{itemize}

A full Gibbs sweep is then the composition
\[
K = K_d \cdots K_2 K_1.
\]
\end{frame}

\begin{frame}{Invariant Distribution: Strategy}
We want to show that \(\pi\) is invariant under Gibbs updates.

\medskip

That means:
\[
\int \pi(dx)\, K(x,A) = \pi(A)
\]
for every measurable set \(A\).

It is enough to prove this first for a single-coordinate update \(K_i\), then invariance for the full sweep follows by composition.
\end{frame}

\begin{frame}{Invariance of a Single-Coordinate Update}
Consider one coordinate \(i\). Then
\[
\int \pi(dx)\,K_i(x,dx')
=
\int \pi(dx_i,dx_{-i})\,\delta_{x_{-i}}(dx'_{-i})\,\pi(dx'_i\mid x_{-i}).
\]

Using the factorization
\[
\pi(dx_i,dx_{-i}) = \pi(dx_i\mid x_{-i})\,\pi_{-i}(dx_{-i}),
\]
we integrate over \(x_i\):
\[
\int \pi(dx_i\mid x_{-i}) = 1.
\]

Hence
\[
\int \pi(dx)\,K_i(x,dx')
=
\pi(dx'_i\mid x'_{-i})\,\pi_{-i}(dx'_{-i})
=
\pi(dx').
\]

So \(\pi\) is invariant for each coordinate update.
\end{frame}

\begin{frame}{Detailed Balance for a Single Update}
In fact, each single-coordinate Gibbs update satisfies detailed balance with respect to \(\pi\):
\[
\pi(dx)\,K_i(x,dx')
=
\pi(dx')\,K_i(x',dx).
\]

Why?

\medskip

Since only coordinate \(i\) changes, we must have \(x_{-i}=x'_{-i}\). Then
\[
\pi(x)\,K_i(x,x')
=
\pi(x_i,x_{-i})\,\pi(x'_i\mid x_{-i}),
\]
and similarly
\[
\pi(x')\,K_i(x',x)
=
\pi(x'_i,x_{-i})\,\pi(x_i\mid x_{-i}).
\]

But using
\[
\pi(x_i,x_{-i})=\pi(x_i\mid x_{-i})\pi(x_{-i}),
\]
both sides become
\[
\pi(x_{-i})\,\pi(x_i\mid x_{-i})\,\pi(x'_i\mid x_{-i}).
\]
\end{frame}

\begin{frame}{Why Detailed Balance Matters}
Detailed balance implies invariance:
\[
\pi(dx)\,K(x,dx')=\pi(dx')\,K(x',dx)
\quad \Longrightarrow \quad
\pi K = \pi.
\]

This gives a simple physical picture:
\begin{quote}
in equilibrium, probability flux from \(x\) to \(x'\) is exactly balanced by the reverse flux.
\end{quote}

For Gibbs sampling:
\begin{itemize}
    \item each coordinate update is reversible,
    \item the full deterministic sweep is invariant,
    \item random-scan Gibbs is also reversible under suitable conditions.
\end{itemize}
\end{frame}

\begin{frame}{Random-Scan Gibbs Sampling}
Instead of updating coordinates in fixed order, one may choose a random coordinate \(I_t\in\{1,\dots,d\}\) at each step, for example uniformly:
\[
\mathbb{P}(I_t=i)=\frac{1}{d}.
\]

Then update
\[
x^{(t+1)}_i \sim \pi(x_i \mid x^{(t)}_{-i}),
\]
while leaving the remaining coordinates unchanged.

The transition kernel becomes
\[
K=\frac{1}{d}\sum_{i=1}^d K_i.
\]

Since each \(K_i\) leaves \(\pi\) invariant, so does \(K\).
\end{frame}

\begin{frame}{Gibbs Sampling as a Special Case of Metropolis--Hastings}
Metropolis--Hastings proposes a move \(x\to x'\) and accepts it with probability
\[
\alpha(x,x')=\min\left(1,\frac{\pi(x')q(x\mid x')}{\pi(x)q(x'\mid x)}\right).
\]

In Gibbs sampling, the proposal is exactly the conditional distribution:
\[
q(x'_i\mid x)=\pi(x'_i\mid x_{-i}).
\]

Then the acceptance ratio becomes
\[
\frac{\pi(x')q(x\mid x')}{\pi(x)q(x'\mid x)}=1,
\]
so every move is accepted.

Thus Gibbs sampling is a rejection-free special case of Metropolis--Hastings.
\end{frame}

\begin{frame}{Connection to the Boltzmann Distribution}
Suppose
\[
\pi(x)=\frac{1}{Z}e^{-\beta E(x)}.
\]

Then the conditional distribution for coordinate \(x_i\) is
\[
\pi(x_i\mid x_{-i})
=
\frac{e^{-\beta E(x_i,x_{-i})}}
{\sum_{x_i'} e^{-\beta E(x_i',x_{-i})}}
\]
in the discrete case.

This is the local Boltzmann law: update one degree of freedom while the environment \(x_{-i}\) is frozen.
\end{frame}

\begin{frame}{Example: Ising Model}
For Ising spins \(s_i\in\{-1,+1\}\), the energy is
\[
E(s)=-\sum_{\langle i,j\rangle} J_{ij}s_i s_j - \sum_i h_i s_i.
\]

The Boltzmann distribution is
\[
\pi(s)=\frac{1}{Z}e^{-\beta E(s)}.
\]

To update spin \(s_i\), only terms involving \(s_i\) matter:
\[
E(s)= - s_i \left(\sum_{j\neq i} J_{ij}s_j + h_i\right) + \text{const}.
\]

Hence
\[
\pi(s_i\mid s_{-i})
\propto
\exp\left[
\beta s_i \left(\sum_{j\neq i} J_{ij}s_j + h_i\right)
\right].
\]
\end{frame}

\begin{frame}{Explicit Conditional for the Ising Spin}
Define the local field
\[
H_i^{\mathrm{loc}}=\sum_{j\neq i} J_{ij}s_j + h_i.
\]

Then
\[
\mathbb{P}(s_i=+1\mid s_{-i})
=
\frac{e^{\beta H_i^{\mathrm{loc}}}}
{e^{\beta H_i^{\mathrm{loc}}}+e^{-\beta H_i^{\mathrm{loc}}}}
=
\frac{1}{1+e^{-2\beta H_i^{\mathrm{loc}}}}.
\]

Similarly,
\[
\mathbb{P}(s_i=-1\mid s_{-i})
=
\frac{1}{1+e^{2\beta H_i^{\mathrm{loc}}}}.
\]

This logistic form is one of the most important explicit examples of Gibbs sampling in physics.
\end{frame}

\begin{frame}{Continuous Example: Bivariate Gaussian}
Let
\[
\begin{pmatrix}
X\\Y
\end{pmatrix}
\sim \mathcal{N}
\left(
\begin{pmatrix}
\mu_X\\\mu_Y
\end{pmatrix},
\begin{pmatrix}
\sigma_X^2 & \rho \sigma_X \sigma_Y \\
\rho \sigma_X \sigma_Y & \sigma_Y^2
\end{pmatrix}
\right).
\]

Then the conditional distributions are Gaussian:
\[
X\mid Y=y \sim \mathcal{N}\left(\mu_X+\rho\frac{\sigma_X}{\sigma_Y}(y-\mu_Y),\;\sigma_X^2(1-\rho^2)\right),
\]
\[
Y\mid X=x \sim \mathcal{N}\left(\mu_Y+\rho\frac{\sigma_Y}{\sigma_X}(x-\mu_X),\;\sigma_Y^2(1-\rho^2)\right).
\]

So Gibbs sampling alternates exact Gaussian draws.
\end{frame}

\begin{frame}{Convergence: What Must Be True?}
Invariance alone is not enough. We also need the Markov chain to converge to \(\pi\).

Typical sufficient conditions are:
\begin{itemize}
    \item irreducibility,
    \item aperiodicity,
    \item positive recurrence.
\end{itemize}

Under such conditions, Gibbs sampling satisfies
\[
\lim_{t\to\infty}\mathbb{P}(X^{(t)}\in A)=\pi(A)
\]
for suitable initial conditions.

So the chain forgets its initial state and approaches equilibrium.
\end{frame}

\begin{frame}{Ergodic Theorem}
If the Gibbs chain is ergodic and \(f\) is integrable, then
\[
\frac{1}{N}\sum_{t=1}^N f(X^{(t)})
\longrightarrow
\mathbb{E}_\pi[f]
\]
almost surely as \(N\to\infty\).

This is the practical reason Gibbs sampling is useful:
we can estimate expectations under \(\pi\) by time averages along the chain.

\medskip

Examples:
\begin{itemize}
    \item magnetization in the Ising model,
    \item posterior means in Bayesian inference,
    \item correlation functions in lattice models.
\end{itemize}
\end{frame}

\begin{frame}{Why Gibbs Sampling Can Be Slow}
Even though the algorithm is simple, convergence may be slow if:
\begin{itemize}
    \item variables are strongly correlated,
    \item the target distribution is multimodal,
    \item local updates move only small distances,
    \item there is critical slowing down near phase transitions.
\end{itemize}

This is especially important in statistical physics, where local Gibbs updates may decorrelate very slowly near criticality.
\end{frame}

\begin{frame}{Blocked Gibbs Sampling}
A standard improvement is to update blocks of variables together.

Instead of sampling one coordinate \(x_i\), sample a block \(x_B\):
\[
x_B^{(t+1)} \sim \pi(x_B \mid x_{-B}^{(t)}).
\]

This can reduce autocorrelation when variables inside the block are strongly coupled.

\medskip

So the Gibbs philosophy remains the same:
\begin{quote}
sample exactly from conditional distributions
\end{quote}
but the conditionals are taken for larger subsets.
\end{frame}

\begin{frame}{Mathematical Summary of the Derivation}
The logic of Gibbs sampling is:
\begin{enumerate}
    \item start from a target distribution \(\pi(x_1,\dots,x_d)\),
    \item identify tractable conditional laws \(\pi(x_i\mid x_{-i})\),
    \item define a Markov chain by sequentially sampling from these conditionals,
    \item prove that each coordinate update preserves \(\pi\),
    \item conclude that \(\pi\) is invariant for the full chain,
    \item under ergodicity conditions, conclude convergence to \(\pi\).
\end{enumerate}
\end{frame}

\begin{frame}{Physical Interpretation}
In statistical physics, Gibbs sampling can be interpreted as:
\begin{itemize}
    \item local equilibration of one degree of freedom,
    \item while the surrounding environment is held fixed,
    \item repeated many times until global equilibrium is reached.
\end{itemize}

This makes it a natural algorithm for:
\begin{itemize}
    \item spin systems,
    \item lattice field models,
    \item Boltzmann machines,
    \item Bayesian graphical models.
\end{itemize}
\end{frame}

\begin{frame}{Final Take-Home Message}
Gibbs sampling is mathematically simple but conceptually deep.

\medskip

It works because:
\begin{itemize}
    \item conditional distributions are exact,
    \item the induced Markov chain preserves the full target distribution,
    \item detailed balance and invariance guarantee correctness,
    \item ergodicity guarantees convergence.
\end{itemize}

So Gibbs sampling is one of the clearest examples of how local probabilistic updates can reconstruct a global equilibrium distribution.
\end{frame}

\end{document}
